\section{5C Lebesgue integration on $R^n$}
\addcontentsline{toc}{section}{5C Lebesgue integration on $R^n$}

3 This should morally be the case, but is surprisingly tricky to prove.
4 Shows that the Lebesgue measure is translation invariant. In fact, we can define the Lebesgue measure on $R^n$ axiomatically as the unique sigma-finite measure that is translation invariant
5 Some version of change of variables for Lebesgue integrals. Axler should really have done the more general version, but that uses determinants and he doesn't want to admit that they're actually useful
6 In other words, the product of complete sigma-algebras is definetely not complete. The Lebesgue measurable subsets are really not very well-behaved. Barry Simon in his 5 comprehensive analysis books says never to use Lebesgue measurable and always use Borel. He calls Lebesgue measurable to be "invented by the devil"
7
12 This is very interesting! Apparently the unit ball takes up less and less space as the dimension gets higher. This hsould remind you of the curse of dimensionality!
13 A general formula for the size of the n-dimensional ball

\begin{exercise}{3}
fill
\end{exercise}
\begin{proof}
fill
\end{proof} 

\begin{exercise}{4}
fill
\end{exercise}
\begin{proof}
fill
\end{proof} 

\begin{exercise}{5}
fill
\end{exercise}
\begin{proof}
fill
\end{proof} 

\begin{exercise}{6}
fill
\end{exercise}
\begin{proof}
fill
\end{proof} 

\begin{exercise}{7}
fill
\end{exercise}
\begin{proof}
fill
\end{proof} 

\begin{exercise}{12}
fill
\end{exercise}
\begin{proof}
fill
\end{proof} 

\begin{exercise}{13}
fill
\end{exercise}
\begin{proof}
fill
\end{proof} 
